\chapter{Краткое описание методов}

\section{Метод исключения Гаусса с выбором главных элементов}

Прямые методы решения систем линейных уравнений позволяют найти точное решение (с точностью до погрешностей машинного округления) за конечное число арифметических операций. Базовым прямым методом является метод последовательного исключения неизвестных Гаусса.

Классический метод Гаусса состоит из двух основных фаз:
\begin{enumerate}
    \item \textbf{Прямой ход}: преобразование исходной системы с помощью эквивалентных элементарных операций над строками расширенной матрицы к эквивалентной системе с верхней треугольной матрицей коэффициентов;
    \item \textbf{Обратный ход}: последовательное нахождение компонент вектора решения, начиная с последнего уравнения (с одним неизвестным) к первому.
\end{enumerate}

\subsection{Проблема устойчивости и выбор главного элемента}

В стандартной схеме Гаусса на $k$-м шаге ведущим (опорным) выбирается диагональный элемент $a_{kk}^{(k)}$. Однако такой подход обладает двумя существенными недостатками:
\begin{enumerate}
    \item если $a_{kk}^{(k)} = 0$, алгоритм аварийно завершается из-за невозможности выполнить деление на ноль;
    \item если $|a_{kk}^{(k)}| \ll 1$ (элемент отличен от нуля, но близок к нему), вычисление множителей вида $m_i = -a_{ik}^{(k)} / a_{kk}^{(k)}$ приводит к делению на малую величину. В результате промежуточные коэффициенты строк многократно возрастают, и при их последующем сложении с другими строками происходит потеря значащих цифр из-за вычитания близких чисел (катастрофическая потеря точности).
\end{enumerate}

Для устранения указанных недостатков применяется \textbf{метод Гаусса с выбором главного элемента}. В реализованном алгоритме используется схема выбора главного элемента по всей активной подматрице (полный выбор):
\begin{equation}
    |a_{p, q}| = \max_{i, j} |a_{ij}|,
    \label{eq:pivot}
\end{equation}
где $p$~--- номер главной строки, $q$~--- номер главного столбца, соответствующего неизвестной $x_q$.

После нахождения главного элемента вычисляются множители для всех остальных строк $i \neq p$:
\begin{equation}
    m_i = -\frac{a_{iq}}{a_{pq}}.
\end{equation}
К каждой $i$-й строке ($i \neq p$) прибавляется $p$-я строка, умноженная на $m_i$:
\begin{equation}
    a'_{ij} = a_{ij} + m_i a_{pj}, \quad b'_i = b_i + m_i b_p.
\end{equation}
В результате столбец $q$ во всех строках, кроме главной строки $p$, обращается в ноль. Далее главная строка $p$ и главный столбец $q$ исключаются из рассмотрения, и процедура рекурсивно повторяется для подматрицы порядка $n-1$.

\subsection{Анализ вырождения при нулевом главном элементе}

В соответствии с замечанием к методу, вычисления нельзя производить, если на каком-либо шаге прямого хода диагональный элемент равен нулю. Перестановкой строк отличный от нуля элемент перемещают на главную диагональ и продолжают вычисления. 

Если же оказывается, что все элементы столбца ниже диагонального (и в текущей подматрице) также равны нулю (с машинной точностью $\varepsilon_{\text{mach}} = 10^{-9}$):
\begin{equation}
    \max_{i, j} |a_{ij}| \le \varepsilon_{\text{mach}},
\end{equation}
то определитель системы равен нулю, и вычисляется характер решения по оставшимся свободным членам $b_i$:
\begin{enumerate}
    \item если хотя бы для одной строки свободный член отличен от нуля ($|b_i| > \varepsilon_{\text{mach}}$), то получается противоречие $0 \cdot x_1 + \ldots + 0 \cdot x_n = b_i \neq 0$. В этом случае система \textbf{не имеет решения};
    \item если для всех оставшихся строк свободные члены равны нулю ($|b_i| \le \varepsilon_{\text{mach}}$), получается тождество $0 = 0$. В этом случае система имеет \textbf{бесконечное множество решений}.
\end{enumerate}

\subsection{Обратный ход метода Гаусса}

Если вырождения нет, то после завершения прямого хода неизвестные восстанавливаются в порядке, обратном их исключению. Для каждой главной строки $p$ с главным столбцом $q$:
\begin{equation}
    x_q = \frac{1}{a_{pq}} \left( b_p - \sum_{j \neq q} a_{pj} x_j \right),
    \label{eq:gauss_back}
\end{equation}
где все $x_j$ ($j \neq q$) уже вычислены на предыдущих этапах обратного хода.

При использовании метода Гаусса производится $\frac{1}{3}n(n^2 + 3n - 1)$ операций умножения и деления и $\frac{1}{6}n(n - 1)(2n + 5)$ вычитаний, то есть всего $\approx \frac{2}{3}n^3$ арифметических операций.

\section{Метод простой итерации}

При большом числе неизвестных (более $10^3$) в системе линейных уравнений применение точных методов становится затруднительным. В таких случаях для решения системы удобнее оказываются итерационные методы.

\subsection{Приведение системы к виду, удобному для итераций}

Метод простой итерации состоит в том, что система уравнений $Ax = b$ преобразуется к виду:
\begin{equation}
    x = Bx + c,
    \label{eq:iter_canonical}
\end{equation}
которая решается методом последовательных приближений. За нулевое приближение принимается вектор свободных членов $x^{(0)} = c$. Далее вычисляются последовательные приближения по формуле:
\begin{equation}
    x^{(k)} = Bx^{(k-1)} + c, \quad k = 1, 2, 3, \ldots
    \label{eq:iter_formula}
\end{equation}
Если последовательность приближений имеет предел, то этот предел является решением системы.

Преобразовать систему $Ax = b$ к виду $x = Bx + c$ можно, выделив диагональные элементы:
\begin{equation}
    x_i = \frac{1}{a_{ii}} \left( b_i - \sum_{j \neq i} a_{ij} x_j \right), \quad i = 1, \ldots, n.
\end{equation}
В матричной форме элементы матрицы $B$ и вектора $c$ определяются выражениями:
\begin{equation}
    b_{ij} =
    \begin{cases}
        0, & i = j; \\
        -\dfrac{a_{ij}}{a_{ii}}, & i \neq j,
    \end{cases}
    \qquad
    c_i = \dfrac{b_i}{a_{ii}}.
    \label{eq:b_and_c}
\end{equation}

\subsection{Условия сходимости и нормы матриц}

Для сходимости метода простой итерации достаточно, чтобы выполнялось условие $\|B\| < 1$ по какой-либо норме матрицы.

Норму матрицы называют согласованной с векторной нормой, если для любого вектора $x$ выполняется условие:
\begin{equation}
    \|Bx\| \le \|B\| \cdot \|x\|.
\end{equation}
Норму матрицы часто вводят через нормы векторов, полагая:
\begin{equation}
    \|B\| = \sup_{x \neq 0} \frac{\|Bx\|}{\|x\|}.
\end{equation}
Такую норму матрицы называют матричной нормой, подчиненной векторной норме.

В соответствии с заданием используются две нормы:
\begin{enumerate}
    \item \textbf{Октаэдрическая норма}: для октаэдрической нормы вектора $\|x\|_1 = \sum_{j=1}^n |x_j|$ подчиненной матричной нормой является максимум сумм модулей элементов по столбцам:
    \begin{equation}
        \|B\|_1 = \max_{1 \le j \le n} \sum_{i=1}^n |b_{ij}|;
        \label{eq:norm_oct}
    \end{equation}
    \item \textbf{Кубическая норма}: для кубической нормы вектора $\|x\|_\infty = \max_{1 \le j \le n} |x_j|$ подчиненной матричной нормой является максимум сумм модулей элементов по строкам:
    \begin{equation}
        \|B\|_\infty = \max_{1 \le i \le n} \sum_{j=1}^n |b_{ij}|.
        \label{eq:norm_cub}
    \end{equation}
\end{enumerate}

\subsection{Обеспечение диагонального преобладания}

Матрица $A$ называется матрицей со \textbf{строгим диагональным преобладанием}, если для каждой ее строки модуль диагонального элемента строго превосходит сумму модулей всех остальных элементов этой строки:
\begin{equation}
    |a_{ii}| > \sum_{j \neq i} |a_{ij}|, \quad \forall i = 1, \ldots, n.
    \label{eq:diag_dom}
\end{equation}

Из соотношений~\eqref{eq:b_and_c} и~\eqref{eq:norm_cub} непосредственно следует:
\begin{equation}
    \|B\|_\infty = \max_{1 \le i \le n} \sum_{j \neq i} \frac{|a_{ij}|}{|a_{ii}|}.
\end{equation}
Следовательно, условие строгого диагонального преобладания~\eqref{eq:diag_dom} является гарантией того, что кубическая норма матрицы перехода строго меньше единицы:
\begin{equation}
    |a_{ii}| > \sum_{j \neq i} |a_{ij}| \iff \sum_{j \neq i} \frac{|a_{ij}|}{|a_{ii}|} < 1 \implies \|B\|_\infty < 1.
\end{equation}

Если исходная матрица системы не удовлетворяет условию диагонального преобладания, ее предварительно преобразуют в эквивалентную систему путем взятия линейных комбинаций строк. В разработанной программе реализован алгоритм перебора коэффициентов линейных комбинаций $w_r \in \{-1, 0, 1\}$:
\begin{equation}
    \operatorname{row}'_i = \sum_{r=1}^n w_r \operatorname{row}_r,
\end{equation}
позволяющий синтезировать новую систему с доминирующей главной диагональю.

\subsection{Критерий окончания итерационного процесса}

В качестве критерия окончания итерационного процесса используется условие малости разности двух последовательных приближений:
\begin{equation}
    \|x^{(k)} - x^{(k-1)}\|_\infty = \max_{1 \le i \le n} |x_i^{(k)} - x_i^{(k-1)}| < \varepsilon,
    \label{eq:stop_crit}
\end{equation}
где $\varepsilon = 10^{-5}$~--- заданная точность вычислений. При выполнении этого условия процесс итераций завершается, а полученный вектор $x^{(k)}$ принимается в качестве искомого решения.
